\chapter{Implementation}
\label{ch:implementation}

This chapter details the concrete software engineering implementation of the Wolverine platform. It documents the core TypeScript analytical backend, database deployment, the WolverineDB cryptographic middleware, the distinct roles of the two user-interface systems, the containerized multi-network deployment, configuration management, and failure recovery mechanisms.

\section{Backend Architecture and Modules}
\label{sec:impl-backend}
The analytical backend resides in \texttt{wolverine-sih/wolverine/} and is implemented in TypeScript (Node.js 20 runtime). The service architecture is structured into six core operational modules:

\begin{itemize}
    \item \textbf{Collector Adapters (\texttt{src/collector/adapters.ts}):} Implements protocol-specific HTTP clients configured with proxy routing to query target onion endpoints via the Tor gateway. Ingested responses are assigned a unique \texttt{captureId} and streamed directly to MinIO.
    \item \textbf{Normalization Parsers (\texttt{src/normalizer/parsers.ts}):} Contains five format-specific parser implementations (\texttt{AtlasParser}, \texttt{BriarParser}, \texttt{CinderParser}, \texttt{DriftParser}, \texttt{EmberParser}). Utilizes Cheerio for server-rendered HTML traversal (Briar Bazaar), \texttt{JSON:API} deserializers (Cinder Exchange), and GraphQL response decoders (Ember Commons).
    \item \textbf{Entity Resolver (\texttt{src/resolver/entity\_resolver.ts}):} Implements the candidate blocking and weighted similarity evaluation algorithms formalized in \Cref{ch:analytical-methodology}.
    \item \textbf{Graph Projector (\texttt{src/graph/projector.ts}):} Maintains the in-memory directed graph using ES6 \texttt{Map} data structures, providing multi-hop BFS traversal and edge-weight filtering.
    \item \textbf{AI/RAG Layer (\texttt{src/ai/rag.ts}):} Wraps local Ollama HTTP endpoints, enforcing the 4-rule input sanitization pipeline and deterministic fallback generation.
    \item \textbf{API Server (\texttt{src/api/server.ts}):} Exposes an Express REST API on port 4000, serving endpoints for collection execution, record inspection, graph queries, entity resolution review, and AI case synthesis.
\end{itemize}

\section{Database and Infrastructure Implementation}
\label{sec:impl-database}
The platform orchestrates six independent database services in Docker containers:
\begin{itemize}
    \item \texttt{postgres-sites}: PostgreSQL 16 image initializing four distinct databases (\texttt{atlas\_market}, \texttt{briar\_bazaar}, \texttt{drift\_forum}, \texttt{ember\_commons}) via initialization SQL scripts.
    \item \texttt{mysql-site-c}: MySQL 8.4 container pre-configured with InnoDB engine for Cinder Exchange.
    \item \texttt{redis}: Redis 7.2 container configured with persistent append-only files (AOF).
    \item \texttt{minio}: MinIO container exposing S3 REST APIs on port 9000 with a dedicated \texttt{captures} bucket.
    \item \texttt{postgres-wolverine}: PostgreSQL 16 container storing analytical canonical records and resolved entity edges.
    \item \texttt{postgres-truth}: PostgreSQL 16 container attached exclusively to the isolated \texttt{wolverine-truth} network with credentials restricted to the evaluation harness.
\end{itemize}

\section{WolverineDB Middleware Realization}
\label{sec:impl-wolverinedb}
\texttt{wolverine-db} is organized across 139 source files in \texttt{wolverine-db/src/}:
\begin{itemize}
    \item \texttt{src/crypto/}: Implements RFC 6962 Merkle tree construction (\texttt{merkle.ts}) and Ed25519 multi-signature validation (\texttt{approval.ts}).
    \item \texttt{src/postgres/} and \texttt{src/adapters/}: Implements Change Data Capture (CDC) triggers and SQLite/PostgreSQL storage drivers.
    \item \texttt{src/bft\_hardening/}: Implements 4-of-5 Byzantine fault tolerance quorum state machines.
    \item \textbf{Verified Simulation Boundaries:} Inter-node network communication is realized via \texttt{DirectMemoryNetworkTransport} (in-process event queues), and EVM state anchoring is implemented via an in-memory registry map (\texttt{src/anchors/evm.ts}).
\end{itemize}

\section{Frontend Systems: Analyst UI vs. Vzeya}
\label{sec:impl-frontend}
A critical finding of the forensic codebase audit is the strict separation between the functional analyst tool and the demonstration UI:

\subsection{Analyst UI (Functional Interface)}
Located in \texttt{wolverine-sih/analyst-ui/}, this is a single-page application built with React 18, Vite, and TailwindCSS. It is containerized and served on port 80/443, actively querying the \texttt{wolverine-api} (port 4000) to render live canonical records, execute interactive graph traversals, and submit sanitized RAG prompts.

\subsection{Vzeya (Cinematic Presentation Mockup)}
Located in \texttt{d:\textbackslash Vault\textbackslash Vzeya/}, Vzeya is a Next.js 15 application utilizing React 19, Three.js, GSAP, and custom WebGL shaders (providing CRT scanline post-processing). Forensic inspection reveals that Vzeya contains no backend API directory (\texttt{src/app/api/} does not exist) and no database connection libraries. All tactical HUD displays, terminal intercepts, and IP network topologies operate entirely over static in-memory data structures for visual demonstration and cinematic storytelling.

\section{Deployment and Network Segmentation}
\label{sec:impl-deployment}
Deployment is configured in \texttt{docker-compose.yml} with three bridge networks:
\begin{itemize}
    \item \texttt{wolverine-dmz}: Bridged network for \texttt{analyst-ui} and \texttt{tor-gateway}.
    \item \texttt{wolverine-internal}: Internal network for source sites, databases, and the analytical engine.
    \item \texttt{wolverine-truth}: Internal-only network (\texttt{internal: true}) with no external gateway for ground-truth isolation.
\end{itemize}

\section{Failure Handling and Graceful Degradation}
\label{sec:impl-failure-handling}
The implementation incorporates specific resilience mechanisms:
\begin{enumerate}
    \item \textbf{LLM Timeout \& Fallback:} If the local Ollama daemon fails or exceeds a 5000\,ms timeout, \texttt{rag.ts} automatically transitions to deterministic heuristic extraction.
    \item \textbf{Tor Gateway Failover:} If the Tor proxy is offline, collector adapters can be toggled to direct internal container hostnames for local development testing.
    \item \textbf{Schema Validation Errors:} Invalid raw payloads are rejected during canonical normalization and flagged with provenance capture hashes for offline debugging.
\end{enumerate}
